\begin{figure}[!t]
  \centering
  \begin{tikzpicture}[
      font=\scriptsize,
      box/.style={draw=figdata!75, fill=white, rounded corners=2pt, align=left, inner sep=3pt,
                  execute at begin node=\hyphenpenalty 10000\relax},
      pbox/.style={box, draw=figprofile!80},
      ubox/.style={box, draw=figudv!85},
      sbox/.style={box, draw=figsim!80},
      lane/.style={rounded corners=3pt, inner sep=4pt},
      plane/.style={lane, draw=figprofile!60, fill=figprofile!7},
      ulane/.style={lane, draw=figudv!65, fill=figudv!8},
      slane/.style={lane, draw=figsim!60, fill=figsim!7},
      heading/.style={font=\scriptsize\bfseries, inner sep=0pt, anchor=south west},
      chunk/.style={align=left, inner sep=0pt, anchor=north west, text width=4.45cm,
                    execute at begin node=\hyphenpenalty 10000\relax},
      digit/.style={font=\tiny, text=black!60, inner sep=0pt},
      on/.style={circle, inner sep=0pt, minimum size=4.2pt, draw=figsim!85, fill=figsim!85},
      off/.style={circle, inner sep=0pt, minimum size=4.2pt, draw=figsim!85, fill=white},
      cell/.style={draw=figsim!60, fill=white},
      link/.style={-{Stealth[length=4.5pt]}, draw=black!65, line width=0.6pt},
      plink/.style={link, draw=figprofile},
      ulink/.style={link, draw=figudv},
      slink/.style={link, draw=figsim},
      note/.style={font=\scriptsize, inner sep=1.5pt},
      pnote/.style={note, text=figprofile!85!black},
      unote/.style={note, text=figudv!85!black},
    ]
    \node[pbox, text width=3.0cm, anchor=north] (per) at (1.65,-0.3)
      {\textbf{Perfil} de até 800 palavras, com itens numerados:\\[1pt]
       \begin{tabular}{@{}l@{\ }l@{}}
         P1, \ldots & posições\\
         C1, \ldots & critérios e valores\\
         A1, \ldots & alinhamentos\\
                    & declarados\\
         F1, \ldots & forma de argumentar
       \end{tabular}};

    \node[chunk] (s0) at (4.95,-0.25) {\textbf{\textit{System prompt}}};
    \node[chunk] (s1) at ($(s0.south west)+(0,-3pt)$)
      {Você simula como \textit{nome}, \textit{cargo sem partido}, se posiciona em audiências
       públicas da Câmara dos Deputados. [\ldots] \textcolor{black!55}{e três instruções}};
    \node[chunk] (s3) at ($(s1.south west)+(0,-4pt)$)
      {PERFIL \textcolor{black!55}{(``Nenhum.'' na condição 0)}\\\textit{itens numerados}};

    \node[chunk] (m0) at ($(s3.south west)+(0,-13pt)$) {\textbf{Mensagem do usuário}};
    \node[chunk] (m1) at ($(m0.south west)+(0,-3pt)$)
      {AUDIÊNCIA DE \textit{data}, SOBRE: \textit{assunto}};
    \node[chunk] (m2) at ($(m1.south west)+(0,-4pt)$)
      {TRECHOS DE FALAS ANTERIORES DE \textit{nome}\\
       {}[T1] \textit{data} $\cdot$ \textit{assunto de origem}\\
       ``\textit{sentença recuperada, com a anterior e a seguinte}''\\
       \ldots\ até [T$k$]};
    \node[chunk] (m3) at ($(m2.south west)+(0,-4pt)$)
      {Uma destas afirmações foi feita por \textit{nome} nesta audiência, que não aparece no
       material desta conversa. Qual é a mais provável?\\
       A) $d_1$\enspace B) $o$\enspace C) $d_2$\enspace D) $d_3$\enspace
       \textcolor{black!55}{(rotação 1)}\\
       Responda só com a letra, mesmo sem certeza.};

    \node[pbox, text width=3.0cm] (ftr) at (1.65,0 |- m2)
      {\textbf{Falas de treino} da pessoa, por audiência e em ordem de data};
    \node[ubox, text width=3.0cm, anchor=north] (aud) at ($(ftr.south)+(0,-0.8)$)
      {\textbf{Resumo estruturado} de uma audiência de validação ou de teste\\[2pt]
       $o$: opinião atribuída à pessoa, ligada a ela pela UDV\\[2pt]
       $d_1$, $d_2$, $d_3$: opiniões de três outros participantes, sorteados};

    \node[digit, anchor=east] at (9.49,0 |- s0) {condição};
    \foreach \row in {s0, m0}
      \foreach \x/\d in {9.62/0, 9.89/1, 10.16/2}
        \node[digit] at (\x,0 |- \row) {\d};
    \foreach \row/\a/\b/\c in {s1/1/1/1, s3/0/1/1, m1/1/1/1, m2/0/0/1, m3/1/1/1}
      \foreach \x/\on in {9.62/\a, 9.89/\b, 10.16/\c} {
        \ifnum\on=1 \node[on] at (\x,0 |- \row) {}; \else \node[off] at (\x,0 |- \row) {}; \fi
      }
    \coordinate (sysr) at (10.24,0 |- s0);
    \coordinate (msgr) at (10.24,0 |- m0);

    \node[align=center, text width=3.3cm, inner sep=0pt, anchor=north] (rtt) at (12.63,-0.3)
      {Opções nas quatro rotações (linhas) de uma ordem sorteada};
    \foreach \l [count=\c] in {A, B, C, D}
      \node[font=\scriptsize\bfseries] at (11.13+0.6*\c,-1.17) {\l};
    \foreach \r/\oc in {1/2, 2/1, 3/4, 4/3} {
      \node[digit] at (11.3,-1.51-0.36*\r+0.36) {\r};
      \foreach \c in {1,...,4} {
        \ifnum\c=\oc
          \draw[cell, fill=figudv!30] (10.85+0.6*\c,-1.35-0.36*\r+0.36) rectangle ++(0.56,-0.32);
        \else
          \draw[cell] (10.85+0.6*\c,-1.35-0.36*\r+0.36) rectangle ++(0.56,-0.32);
        \fi
      }
    }
    \foreach \r/\l in {1/{d_1,o,d_2,d_3}, 2/{o,d_2,d_3,d_1}, 3/{d_2,d_3,d_1,o}, 4/{d_3,d_1,o,d_2}}
      \foreach \t [count=\c] in \l
        \node at (11.13+0.6*\c,-1.51-0.36*\r+0.36) {$\t$};
    \coordinate (grid) at (12.63,-2.75);
    \coordinate (gridw) at (11.2,-1.2);

    \node[sbox, text width=2.7cm, align=center, draw=figsim!90, fill=figsim!22, anchor=north]
      (mod) at (12.63,-3.1)
      {\textbf{Modelo simulador}\\uma chamada por rotação em cada condição};
    \node[sbox, text width=3.0cm, below=0.3cm of mod] (prb)
      {Probabilidade de A, B, C e D no primeiro token, renormalizada entre as quatro em cada
       rotação};
    \node[sbox, text width=3.0cm, below=0.3cm of prb] (avg)
      {Média de cada opção nas quatro rotações; a resposta é a de maior probabilidade};
    \node[sbox, text width=3.0cm, below=0.3cm of avg] (met)
      {Por condição: acerto e probabilidade média de $o$};

    \begin{scope}[on background layer]
      \node[sbox, fit=(s0)(s1)(s3)(sysr)] (sys) {};
      \node[sbox, fit=(m0)(m1)(m2)(m3)(msgr)] (msg) {};
      \node[slane, fit=(sys)(msg)] (conv) {};
      \node[plane, fit=(per)(ftr)] (bperf) {};
      \node[ulane, fit=(aud)] (baud) {};
      \node[slane, fit=(rtt)(mod)(met)(gridw)] (read) {};
    \end{scope}
    \node[heading, text=figprofile!85!black] at ($(bperf.north west)+(0,2pt)$) {Perfil e trechos};
    \node[heading, text=figudv!85!black] at ($(baud.north west)+(0,2pt)$) {Pergunta};
    \node[heading, text=figsim!85!black] at ($(conv.north west)+(0,2pt)$)
      {Conversa com o modelo simulador};
    \node[heading, text=figsim!85!black] at ($(read.north west)+(0,2pt)$) {Leitura da resposta};

    \draw[plink] (ftr) -- node[pnote, right] {LLM gerador} (per);
    \draw[plink] (per.east |- s3) -- (sys.west |- s3);
    \draw[plink] (ftr.east |- m2) -- node[pnote, above] {recuperação} (msg.west |- m2);
    \draw[ulink] (aud.east |- m3) -- node[unote, above] {opções} (msg.west |- m3);
    \draw[slink] (conv.east |- mod) -- (mod.west);
    \draw[slink] (grid) -- (mod.north);
    \draw[slink] (mod) -- (prb);
    \draw[slink] (prb) -- (avg);
    \draw[slink] (avg) -- (met);
  \end{tikzpicture}
  \caption{Pergunta de múltipla escolha da simulação. O perfil e os trechos vêm só das falas de
  treino da pessoa, e a data, o assunto e as opções, de uma audiência de validação ou de teste. Os
  pontos indicam as condições em que cada parte da conversa é incluída (ponto vazio: parte ausente),
  e os campos em itálico são preenchidos a cada chamada (texto completo no
  Apêndice~\ref{ap:prompt-simulacao}). A opinião $o$ ocupa uma posição diferente em cada rotação.}
  \label{fig:simulacao}
\end{figure}
